\documentclass[11pt]{article}
\usepackage[margin=1in]{geometry}
\usepackage{times}
\usepackage{enumitem}
\linespread{1.1}

\begin{document}

\section*{Paper Summaries}

\subsection*{On-device sorting on commodity / programmable hardware}

\subsubsection*{[FPGA but seemingly very relevant] Luan et al.\ 2017 --- Compact standalone platform with real-time on-node spike sorting (bioRxiv 186627)}
Luan, Williams, Constandinou et al.\ built and in-vivo validated a compact, head-mountable standalone platform that performs real-time on-node template-matching spike sorting entirely on off-the-shelf parts: an Intan RHD2132 analog front end feeding a low-power Microsemi IGLOO nano FPGA plus an NXP K64F microcontroller that logs to a micro-SD card. It runs a two-phase scheme---a 2--5 min tethered calibration that clusters raw data offline (Wave\_Clus) into up to four templates per channel, then untethered online detection, peak alignment, and Norm-1 template matching. It handles 32 channels at $\sim$15.8 kSps at $\sim$62--99 mW (the FPGA sorting is $<$5.8\% of power), and was validated live in a freely-moving rhesus macaque (two 24-hour M1 sessions, 21 units), achieving $\sim$823$\times$ data reduction. It is explicitly framed around cutting the wireless bandwidth bottleneck (it even notes an nRF52832 BLE link could stream 79 channels of spike events). Deployed FPGA/MCU system, not an ASIC; logs locally rather than transmitting wirelessly.

\subsubsection*{[FPGA + algorithm] Wu et al.\ 2023 --- Real-time unsupervised opt-OSort on FPGA (J.\ Neural Eng.\ 20 066015)}
Wu et al.\ implement ``opt-OSort,'' an optimized unsupervised OSort streaming spike sorter, on a Xilinx ZYNQ XC7Z020 FPGA (100 MHz) for real-time single-channel sorting. The key change replaces OSort's Euclidean-distance criterion with a Pearson correlation coefficient (amplitude/noise-robust), letting the algorithm converge into a fixed four cluster slots and eliminating OSort's uncontrolled transient-cluster proliferation and iterative merging; two ``loop-checking'' validation loops reject outliers and track electrode drift. Validated on synthetic (OSort, Wave\_Clus) and pre-recorded in-vivo (hc-1 hippocampal, freely-moving-mouse olfactory-bulb) datasets, it reaches accuracy comparable to offline OSort ($\geq$87\%, 79.8\% average, $\sim$92--96\% agreement with offline sorters). Hardware results: 0.68 $\mu$s latency ($\sim$1.47M spikes/s) and 256 bytes/channel. Fabricated/synthesized real-time hardware, unsupervised, single-channel (multiplexable); assumes spikes pre-detected and has no wireless link.

\subsection*{Wireless COTS headstages}

\subsubsection*{[FPGA + optogenetics...]Gagnon-Turcotte et al.\ 2017 --- Wireless optogenetic headstage, on-device detection + compression (IEEE TBioCAS 11(1))}
A COTS wireless optogenetic headstage with 32 recording + 32 optical-stimulation channels that performs spike detection and wavelet data compression in situ in real time. The front end is a commercial Intan RHD2132 (16-bit, sampled here at 20 kSps); all processing runs on a low-power Xilinx Spartan-6 FPGA (Microblaze softcore + custom VHDL: adaptive-threshold detector, four-level Symmlet-2 DWT compression). Detected/compressed AP waveforms are transmitted via a commercial Nordic nRF24L01+ 2.4-GHz transceiver (measured 1.4 Mbps) to a base station 2 m away, with compression up to 508$\times$. Crucially it does NOT sort on-device---PCA + k-means sorting is done offline at the receiver. Measured up to 100\% detection at 15 dB SNR (97.28\% at 5 dB), 31.9 mW/channel; 2.8 g rigid-flex PCB; validated in vivo in a ChR2 transgenic mouse. The single closest COTS prior art on the shared Intan + Nordic axis, but it detects and compresses rather than classifies.


\subsubsection*{[Full wireless streaming] Tasneem et al.\ 2023 --- Duty-cycled wireless optogenetic headstage (IEEE Sensors J.\ 23(21))}
A COTS, 3D-printed ($<$1 g) duty-cycled wireless optogenetic headstage for freely-moving rats combining 16-channel electrophysiology with $\mu$LED stimulation. The front end is a commercial Intan RHD2132; wireless transmission uses a TI CC430 SoC (MSP430 + sub-1-GHz CC1101) at 868 MHz, 500 kBaud, round-robin over 16 channels to a receiver $\sim$2 m away, reconstructed in a LabVIEW GUI. Crucially, both spike detection (continuous wavelet transform) and spike sorting (PCA + k-means) are performed entirely offline in LabVIEW after wireless transmission of the raw neural signal---not on-device. Validated in vivo by recording from the VTA of a freely-moving rat. It is essentially the raw-streaming, sort-offline anti-pattern the thesis aims to beat; useful as a contrasting commodity-hardware baseline.

\subsection*{On-chip spike-sorting ASICs (custom silicon)}

\subsubsection*{Xu et al.\ 2019 --- Unsupervised real-time spike-sorting chip for hippocampal prosthesis (J.\ Neurosci.\ Methods 311)}
A fabricated-and-measured 16-channel spike-sorting chip (40-nm CMOS, 0.0175 mm$^2$/channel, 1.1 V, 8 MHz) for a hippocampal cognitive prosthesis. It splits into a shared training module (dual-vertex-threshold detection + hardware-optimized Osort clustering) and 16 independent real-time classifiers using an optimized Bayes-optimal template-matching algorithm (PBOTM) that skips convolution on noise. Measured (via FPGA-streamed data): 98.3\% detection, 93.4\% classification accuracy, 19 $\mu$W/channel; it handles partially overlapping spikes. Validated on datasets only (synthetic Quiroga, real CRCNS hc-1), not a live animal, with no wireless transmitter. A custom-silicon prior solution to the on-device classification half.

\subsubsection*{Shaeri \& Sodagar 2020 --- Framework for on-implant spike sorting via salient feature selection (Nature Communications 11:3278)}
An on-implant online spike-sorting framework that offloads training to an external ``shadow spike sorter'' and keeps a lightweight engine on the implant, using dynamic salient feature selection followed by multi-label window discrimination. It is a layout/synthesis study, not fabricated silicon: a synthesized 512-channel implementation in 130-nm CMOS occupying 3.36 mm$^2$, 1869 transistors/channel, dissipating 905.9 $\mu$W. Validated on a public dataset only (macaque V1 Utah-array), reaching $\sim$88\% classification accuracy with $\sim$5$\times$ less memory than comparable methods. No fabricated chip and no RF transmitter is built (wireless is assumed conceptual). A high-profile prior on-device classification method, but simulation/synthesis only.

\subsubsection*{Chen et al.\ 2026 --- 128-channel level-crossing digitization + spike-sorting SoC (IEEE TBioCAS)}
A fabricated-and-measured 22-nm FDSOI CMOS SoC integrating 128 level-crossing ADCs with a full digital spike-sorting back-end (LC-based detection, spatiotemporal feature extraction, Geo-Osort clustering); the analog amplifier front end is NOT on-chip. Novelties include LC-pulse-stream central detection, spatial spike realignment, and a compact 8-feature spatiotemporal extractor. Measured: 1.2 $\mu$W/channel (ADC) + 1.09 $\mu$W/channel (sorting core); accuracy 98.37\%/93.52\%/72.39\%/87.17\% across Neuropixels/MEArec/Quiroga sets; $\sim$142.6 $\mu$s latency. Validated on synthetic/replayed datasets, not live animals; no wireless link. A custom-silicon academic competitor, useful for accuracy/power benchmarking.

\subsubsection*{Jiang et al.\ 2026 --- 1024-channel 583 nW/ch spike-sorting SoC (IEEE TCAS-I 73(8))}
A fabricated 28-nm CMOS 1024-channel spike-sorting SoC whose contribution is memory-centric: a sparsity-aware detection scratchpad, a dual-voltage 5T-SRAM buffer, and hybrid on-chip-SRAM/off-chip-FeRAM template memory supporting 16K templates, managed by a firing-rate-based Osort. Measured 583 nW/channel, 11.7 $\mu$s latency; accuracy 91.7\%/88.5\%/84.6\% across Quiroga/Neuropixels/Mouse-MEArec plus 93.7\%/85.1\% on replayed in-vivo data. No integrated analog front end and no wireless transmitter. A custom-silicon competitor demonstrating on-chip classification at very high channel count.

\subsubsection*{Lukito et al.\ 2026 --- Spike-sorting SoC with analog-CIM autoencoder features (IEEE JSSC 61(4))}
A fabricated 65-nm CMOS SoC that performs full spike sorting ($\Delta$-based detection, binary analog computing-in-memory autoencoder feature extraction, L1 K-means clustering) entirely on-chip, and---unlike the other ASICs---includes an 8-channel analog front end (LNA + $\Delta$-subtraction), relocating detection and feature extraction into the analog domain to avoid ADC bottlenecks. Measured 94.54\% Quiroga classification accuracy (highest reported on-chip at the time), 5.11 $\mu$s latency, 3.11 $\mu$W/channel classifier + 5.13 $\mu$W/channel AFE; the $\Delta$-spike approach cuts per-channel data rate 48.8$\times$. Validated on replayed Quiroga + simulated Neuropixels data; no wireless link. The custom-silicon competitor that integrates the AFE.

\subsection*{Algorithm blueprint, benchmarking, and reviews}

\subsubsection*{Navajas \& Quiroga et al.\ 2014 --- Minimum requirements for real-time on-chip spike sorting (J.\ Neurosci.\ Methods 230)}
A computational modeling study (not fabricated hardware) that determines the minimum signal and processing specifications for accurate real-time on-chip spike sorting, explicitly motivated by the wireless-bandwidth bottleneck. It proposes the two-stage hybrid architecture: stage 1 detects spikes on-chip and streams them to an external computer that performs the heavy clustering (Wave\_Clus) to build mean templates, which are sent back; stage 2 then classifies each newly detected spike on-chip against those templates, so only sorted spike identities (not raw data) are transmitted. Using 90 realistic synthetic recordings plus real human medial-temporal-lobe data, it finds sampling can drop to 7 kHz and 10-bit resolution ($\sim$6.4$\times$ reduction) with only $\sim$20\% loss, that squared Euclidean distance is the best performance-vs-complexity metric, that peak alignment is critical, and that a $\sim$14-sample window suffices; it estimates $\sim$2.24 mW for 100 channels ($\sim$20$\times$ below raw transmission). No silicon, but the algorithmic blueprint for the two-phase supervised approach.

\subsubsection*{Zhang et al.\ 2023 --- Survey: spike-sorting algorithms and their hardware implementation (J.\ Neural Eng.\ 20 021001)}
A comprehensive topical review surveying spike-sorting algorithms jointly with their hardware implementations, framed around hardware-algorithm co-optimization for resource-constrained implantable systems. It covers the conventional detect/feature/cluster pipeline, template matching (OSort, KiloSort), deep neural networks, and spiking neural networks, and contrasts off-site (CPU/GPU) with on-site (FPGA, ASIC, memristive in-memory) processing via power/latency/area comparison tables. It concludes there is no single best implementation, that co-design is essential, and discusses drift, overlapping spikes, and the high-density data-rate explosion (Neuropixels $\sim$1 GB/min). Background and design-space map, not itself a prior solution.

\subsubsection*{Magland et al.\ 2020 --- SpikeForest reproducible ground-truth sorter validation (eLife 9:e55167)}
An open-source pipeline and interactive website that benchmarks ten automated spike sorters against a large curated ground-truth database ($\sim$650 recordings, $\sim$35{,}000 units) spanning paired intra/extracellular, simulated, and human-curated data, with accuracy defined per unit at a 1 ms match tolerance. The central conclusion is that no single sorter is best across all conditions (IronClust, KiloSort2, MountainSort4, SpyKING CIRCUS lead in different regimes). A methods/benchmarking resource that supports algorithm selection and the ground-truth-plus-tolerance-window validation methodology, but addresses offline benchmarking on server clusters---not on-device or wireless.

\subsection*{Recording-probe technology (background)}

\subsubsection*{Steinmetz et al.\ 2021 --- Neuropixels 2.0 (Science 372 eabf4588)}
Introduces the Neuropixels 2.0 CMOS probe: 384 simultaneously recordable channels (14-bit), single-shank (1280 sites) or four-shank (5120 sites) at 15 $\mu$m spacing, two probes + headstage $\sim$1.1 g for chronic mouse implants, 30 kHz AP band, $\sim$8.2 $\mu$V RMS noise, streaming 161.3 Mbit/s over one cable; plus a post-hoc motion-correction algorithm and chronic stability over months. A fabricated recording probe defining the state-of-the-art high-channel-count data-rate problem; background/motivation, not an on-device-classification solution.

\subsubsection*{Jun et al.\ 2017 --- Neuropixels 1.0, fully integrated silicon probes (Nature 551)}
The original Neuropixels probe: a fully integrated 130-nm CMOS probe combining 960 recording sites (384 simultaneous channels) with on-chip amplification, multiplexing, and 10-bit digitization on a single 250-mg chip, AP band 0.3--10 kHz sampled at 30 kHz, $\sim$5.1 $\mu$V RMS noise; $>$700 neurons recorded simultaneously with two probes, including grid cells. Foundational integrated-probe engineering establishing the 384-channel stream that motivates the bandwidth problem; background/motivation, not a prior solution.

\subsection*{Foundational neuroscience and physiology (background)}

\subsubsection*{Buzs\'aki, Anastassiou \& Koch 2012 --- The origin of extracellular fields and currents (Nat.\ Rev.\ Neurosci.\ 13)}
A review of the biophysical origins of extracellular fields (EEG, ECoG, LFP, spikes). It explains that the extracellular potential is the superposition of all transmembrane currents, that synaptic currents dominate the LFP, and that fast Na$^+$ action potentials ($<$2 ms) generate the strongest, shortest-duration currents detectable as ``unit''/spike activity, with $>$100 Hz power serving as a proxy for local spiking output. Foundational physiology justifying the spike band and waveform used to generate realistic simulated spikes; background, not a prior solution.

\subsubsection*{Hodgkin \& Huxley 1952 --- Quantitative description of membrane current (J.\ Physiol.\ 117)}
The culminating paper of the Hodgkin-Huxley series, which distilled squid-giant-axon voltage-clamp data into a complete conductance-based model of the action potential (a membrane capacitance in parallel with Na$^+$, K$^+$, and leak branches, with gating variables $g_K=\bar g_K n^4$ and $g_{Na}=\bar g_{Na} m^3 h$ governed by first-order kinetics). Numerically integrating this system reconstructed the AP's form, threshold, refractory period, and conduction velocity with no fitted free parameters. The canonical mathematical model of the action potential and the basis for essentially all conductance-based simulation---directly citable for how the thesis's simulated action potentials are generated; background, not a prior solution.

\subsubsection*{Seyfarth 2006 --- Julius Bernstein: pioneer neurobiologist and biophysicist (Biol.\ Cybern.\ 94)}
A biographical review profiling Julius Bernstein's two landmark contributions: the 1868 differential rheotome (an electromechanical sampling instrument enabling the first quantitative reconstruction of the action potential's $\sim$1 ms time course, measured at $\sim$28.7 m/s conduction velocity) and the 1902 Membrane Theory of electrical potentials. Historical/biographical background supporting the Bernstein citation used in the proposal's Background; not a prior solution.

\subsubsection*{Schuetze 1983 --- The discovery of the action potential (Trends in Neurosciences 6)}
A historical Perspectives article tracing how the action potential was first recorded: Galvani/Volta's animal-electricity debate, du Bois-Reymond's ``negative variation,'' Helmholtz's conduction-velocity measurement, and Bernstein's differential rheotome producing the first published plot of the action potential's time course. Pure historical background, largely redundant with Seyfarth 2006; not a prior solution.

\subsection*{Data-link / telemetry (adjacent sub-problem)}

\subsubsection*{Russo et al.\ 2026 --- Event-based SerDes telemetry network for distributed BCIs (Neuromorph.\ Comput.\ Eng.\ 6 014001)}
A deterministic-latency serializer-deserializer (SerDes) telemetry network moving high-channel-count neural data from multiple implant sensors to a central hub over a wired LVDS star network. It uses event-based send-on-delta sampling via level-crossing ADCs, a ternary packet-based address-event-representation protocol, and crystal-free Manchester coding; the fabricated 65-nm serializer + FPGA deserializer achieve $>$7--11$\times$ compression and $<$6 $\mu$Vrms reconstruction error. It transports full/compressed waveforms and does NOT sort or classify spikes, and the wireless hop is cited rather than built. A wired data-link/compression layer solving a different sub-problem than the thesis; contextual at most.

% ================================================================
\section*{Second Review Round --- Additional Prior Art}

\subsection*{Wireless COTS headstages (continued)}

\subsubsection*{Bilodeau et al.\ 2026 --- Wireless closed-loop headstage with on-device SVM decoding (bioRxiv 2026.03.25.714179)}
A bidirectional wireless optogenetic headstage that performs on-device feature extraction and machine-learning decoding entirely on commodity parts: an Intan RHD2132 analog front end feeds a Xilinx Spartan-6 FPGA, and a Nordic nRF24L01+ 2.4-GHz transceiver ($\sim$2 Mbps) closes the link. The FPGA runs on-device threshold detection, firing-rate binning, and on-chip PCA feeding a nonlinear-SVM decoder, so what is transmitted is compressed spike events plus the decoded behavioral/intention output rather than raw waveforms; spike sorting is done offline only for visualization. It drives closed-loop optogenetic stimulation and was validated in vivo in freely-moving rats. This is the single closest wireless prior art to the thesis on the ``on-device ML + commodity Intan/Nordic + transmit-decisions-not-waveforms'' axis, but it decodes \emph{population} activity to a behavioral variable rather than sorting to per-unit identities, and uses a custom 2.4-GHz link rather than BLE---so the specific gap (on-device single-unit classification, wireless identity transmission, commodity BLE SoC) remains open.

\subsubsection*{Yoon et al.\ 2022 --- Bi-directional wireless drug-delivery/electrophysiology BLE probe (Nature Communications 13:5521)}
A fabricated, head-mounted (4.6 g, $\sim$2 cm$^3$) wireless neural probe that monolithically integrates a 16-electrode silicon shank, microfluidic drug channels, an electrolytic pump, and a bi-directional \emph{Bluetooth Low Energy} module, battery-powered for $\sim$1 h. Neural signals are amplified by a commercial Intan RHD2132, processed by a TI MSP430 MCU, and transmitted via a Panasonic PAN1326B Bluetooth transceiver; to overcome BLE's $<$1 Mbps ceiling the MCU runs real-time threshold-crossing spike detection ($\sim$$-50$ $\mu$V) and transmits only detected spike events plus timing rather than raw waveforms, enabling real-time 16-channel monitoring (74 mW for neural transmission). It was validated in vivo in live, freely-moving, socially-interacting mice with simultaneous wireless drug delivery. This is the closest \textbf{commercial-BLE, commodity-hardware} wireless system to the thesis and shares the exact bandwidth-reduction philosophy, differing only in that it stops at threshold detection rather than on-device classification.

\subsection*{On-device classification --- synthesis / FPGA (no fabricated wireless system)}

\subsubsection*{Siddiqi \& Strydis 2024 --- On-device neural-network spike classification with local non-volatile storage (arXiv 2311.04808)}
An on-device classification design that discriminates cerebellar Purkinje simple vs.\ complex spikes with a compact quantized multilayer-perceptron neural network, targeting an autonomous implant that stores results locally in spin-transfer-torque (STT-RAM) rather than transmitting them. It is a synthesis/layout study (45-nm), not fabricated silicon, validated on datasets only, with no analog front end and no wireless link. It is the closest prior art on the ``\textbf{neural-network classification on-device}'' axis specifically, but its no-wireless, local-storage architecture is the opposite of the thesis's wireless-transmission goal.

\subsubsection*{Han et al.\ 2025 --- ``L-Sort'' localization-based on-device sorting (arXiv 2501.17885)}
``L-Sort'' is a real-time spike sorter that replaces waveform template matching with median-of-median detection followed by \emph{localization-based} clustering---estimating each spike's physical position from its graded amplitude across nearby electrodes (a spatiotemporal, dense-probe feature) and clustering in that space. It is implemented on FPGA and additionally characterized via a post-place-and-route-simulated 180-nm ASIC (no tape-out), validated on datasets, with no wireless link. Relevant as an on-device sorting method that leans on the multi-channel \emph{spatial} feature space the thesis's design section discusses, but it is not a fabricated or wireless system.

\subsubsection*{Valencia \& Alimohammad 2019 --- Real-time parallel-OSort sorting system (IEEE TBioCAS)}
A real-time spike-sorting system built around a parallelized OSort clustering engine, prototyped on FPGA and projected to an estimated 32-nm ASIC. It performs unsupervised on-device clustering with reported throughput/latency figures suitable for many-channel operation, validated on datasets only, with no analog front end and no wireless transmitter. A custom/estimated-silicon on-device unsupervised-sorting competitor, without the commodity-hardware or wireless dimensions.

\subsection*{On-chip spike-sorting ASICs --- custom silicon (continued)}

\subsubsection*{Jiang et al.\ 2020 --- 96-channel 6.3 nW/ch neural spike processor for BCI (arXiv 2009.05210)}
A fabricated-and-measured 96-channel neural spike processor in 180-nm CMOS (1.86 mm$^2$) implementing a full on-implant pipeline: absolute-threshold detection, a novel supervised constrained-Bayesian decision-tree sorter (up to 4 clusters/channel via a 3-bit CAM, 32 detectors sharing one sorter), and partial motor-intention decoding via an Ensemble-Observation Kalman Filter. Its headline is extreme efficiency---6.3 nW/channel (0.61 $\mu$W total at 0.32 V, $\sim$21$\times$ below prior art)---driving cascaded data reduction from 240 kbps/channel down to only $\sim$320 bps total, so it would transmit sorted cluster identities and decoded intention rather than raw signals. Validated offline only (89--99\% sorting accuracy on the Wang mouse dataset; 28\% lower decoding error than a standard Kalman filter on the DREAM Rhesus set); it contains no analog front end and no radio. A custom-silicon prior solution to the on-device-classification-plus-data-reduction half, but not commodity hardware and not wireless.

\subsubsection*{Guo et al.\ 2024 --- 68-channel 22-nm FDSOI processing SoC with on-chip sorting and compression (Front.\ Neurosci.\ 18:1432750)}
A fabricated-and-measured programmable 68-channel neural-processing SoC (9 mm$^2$, 22-nm FDSOI) integrating delta-sigma recording front ends (0.41 $\mu$W/channel), two-stage on-chip detection (adaptive threshold + nonlinear energy operator), intra-/cross-channel compression engines, and a RISC-V core plus MAC accelerator that runs on-chip training and inference for software-defined spike sorting (PCA/adaptive-filter features + K-means/DBSCAN). Measured 91.5\%/94.1\% sorting accuracy on Quiroga, $\sim$91\% near-lossless AP compression (cutting a 1000-channel system's transmission power $\sim$66--79\%), 1.09 $\mu$J/spike. Validated on replayed datasets and in-vitro cultures; crucially it has \emph{no integrated radio}---off-chip transmission is SPI to an external STM32. A high-integration custom-silicon competitor for on-device sorting + compression, but no BLE/wireless and dataset-only validation.

\subsubsection*{Akhoundi et al.\ 2025 --- 1024-channel 40-nm event-driven detection + SOM clustering ASIC (NIH nihms-2121528)}
A fabricated 1024-channel neural-signal processor in 40-nm CMOS performing event-driven spike detection followed by self-organizing-map (SOM) spatial clustering on-chip, at 74 nW/channel, achieving $>$1000$\times$ data reduction. Validated on datasets only, with no analog front end integrated in the processing path and no wireless transmitter. A very-high-channel-count fabricated custom-silicon competitor to the on-device sorting function, without the commodity-hardware or wireless axes.

% ================================================================
\section*{Project the Prior Art Must Be Measured Against}

What I'm trying to build is a working prototype of in-situ classification of simulated neural action potentials for wireless transmission of domain-specific neural information to solve the problem of the Bluetooth Low Energy communication bottleneck to achieve wireless behavioral neuroscientific studies with high channel count. So obviously the wireless communication protocol has to be a part of the system. I will be doing this with off-the-shelf analog front end and SoC for computation, so the novelty I suppose is the algorithmic implementation on that specific hardware and a turnkey solution with however high fidelity I can achieve on simulated data for system neuroscientists to test with their probes on their animals.

% ================================================================
\section*{Most Relevant Prior Art}

The most relevant papers are \textbf{Luan et al.\ 2017}, \textbf{Bilodeau et al.\ 2026}, \textbf{Gagnon-Turcotte et al.\ 2017}, \textbf{Yoon et al.\ 2022}, \textbf{Navajas \& Quiroga et al.\ 2014}, and \textbf{Wu et al.\ 2023}. Together they span the exact axes of the proposed system: on-device real-time sorting on commodity/programmable hardware (Luan, Wu), on-device ML with wireless transmission of decisions rather than waveforms on off-the-shelf Intan + Nordic parts (Bilodeau, Gagnon-Turcotte), the shared \emph{commercial-BLE} bandwidth-reduction philosophy on commodity hardware (Yoon), and the two-stage template-based algorithm that transmits only spike identities (Navajas). Read closely, they sharpen rather than close the gap the thesis fills: \emph{no prior system combines on-device \textbf{classification} (sorting to per-unit identity, not mere detection or population decoding) with \textbf{wireless} transmission of those identities over a \textbf{commodity BLE SoC}}. Luan sorts on-device on commodity hardware but logs to an SD card rather than transmitting; Bilodeau runs on-device ML and transmits wirelessly on commodity Intan + Nordic hardware but decodes \emph{population} activity to a behavioral variable over a custom 2.4-GHz link, not sorted identities over BLE; Yoon transmits over \emph{actual BLE} on commodity Intan hardware but stops at threshold detection; Gagnon-Turcotte transmits wirelessly on Intan + Nordic hardware but only detects and compresses, sorting offline; Tasneem streams raw data and sorts offline; Siddiqi \& Strydis classify on-device with a neural network but store locally with no wireless; and the custom-silicon ASICs (Chen, Jiang, Lukito, Xu, Guo, Akhoundi) classify on-chip but are not commodity hardware and carry no wireless link. The specific combination---on-device single-unit classification, wireless identity transmission, commodity BLE SoC, validated on simulated data---remains unclaimed.

% ================================================================
\section*{Papers to Discard}

\begin{itemize}[nosep]
\item \textbf{Russo et al.\ 2026 (SerDes telemetry)} --- solves a different sub-problem (wired inter-implant serialization/compression, no sorting, not BLE). Cut from Existing Solutions.
\item \textbf{Schuetze 1983} --- largely redundant with Seyfarth 2006 for the Bernstein/action-potential-discovery history; keep one (Seyfarth is the fuller account) and cut this one.
\end{itemize}

\noindent\textbf{Not prior solutions --- reallocate, do not delete.} The following are needed elsewhere in the proposal (Background/Motivation or Methods) but are not ``existing solutions'' and should not clutter that section: \textbf{Jun et al.\ 2017} and \textbf{Steinmetz et al.\ 2021} (Neuropixels 1.0/2.0) and \textbf{Buzs\'aki et al.\ 2012} and \textbf{Hodgkin \& Huxley 1952} and \textbf{Seyfarth 2006} belong in Background/Motivation; \textbf{Zhang et al.\ 2023} (survey) and \textbf{Magland et al.\ 2020} (SpikeForest) belong in the methods/design-space and benchmarking discussion.

% ================================================================
\section*{Prior Solutions, Ranked by Relevance}

\begin{enumerate}[nosep]
\item \textbf{Luan et al.\ 2017} --- on-device real-time sorting (classification) on commodity Intan + FPGA/MCU, freely-moving animal, bandwidth-motivated; missing only the wireless link (logs to SD) and a BLE SoC. Closest overall.
\item \textbf{Bilodeau et al.\ 2026} --- on-device feature extraction + SVM decoding \emph{with} wireless transmission of the decisions, on commodity Intan + Nordic + FPGA, in-vivo closed-loop; differs by decoding \emph{population} activity over a custom 2.4-GHz link rather than sorting to per-unit identities over BLE. Closest wireless prior art.
\item \textbf{Gagnon-Turcotte et al.\ 2017} --- the shared Intan + Nordic COTS wireless philosophy, on-device processing, in-vivo; differs by detecting/compressing and sorting offline rather than classifying on-device.
\item \textbf{Yoon et al.\ 2022} --- commodity Intan + \emph{actual BLE} SoC, on-device threshold detection transmitting spike events under BLE's ceiling, in-vivo freely-moving mice; the closest system on the commercial-BLE axis, but detection-only, no classification.
\item \textbf{Navajas \& Quiroga et al.\ 2014} --- the exact two-stage ``build templates offline, classify on-device, transmit identities'' algorithm; differs by being simulation/analysis only, no hardware.
\item \textbf{Wu et al.\ 2023} --- real-time on-device \emph{unsupervised} sorting on commodity FPGA with tiny latency/memory; differs by having no wireless link and using an FPGA rather than a BLE SoC.
\item \textbf{Siddiqi \& Strydis 2024} --- on-device \emph{neural-network} classification (quantized MLP), the closest match on the NN-on-device axis; differs by storing results locally in STT-RAM with no wireless, synthesis-only.
\item \textbf{Gagnon-Turcotte et al.\ 2015} --- 2-channel commodity Nordic + MCU on-device detection; the methodological predecessor, detection-only.
\item \textbf{Tasneem et al.\ 2023} --- commodity Intan + wireless + in-vivo freely-moving rat, but streams raw and sorts offline; the anti-pattern the thesis beats, valuable as a direct baseline.
\item \textbf{Xu et al.\ 2019} --- fabricated on-chip classification ASIC (template matching + Osort); the on-device classification half, but custom silicon and no wireless.
\item \textbf{Jiang et al.\ 2020} --- fabricated 96-channel on-chip detection + supervised decision-tree sorter + intention decoding at 6.3 nW/ch; custom-silicon on-device-classification + data-reduction prior art, dataset-only, no AFE, no wireless.
\item \textbf{Guo et al.\ 2024} --- fabricated 68-channel 22-nm SoC with on-chip sorting + compression and a RISC-V core; custom-silicon competitor, but SPI to an external MCU (no BLE), dataset/in-vitro only.
\item \textbf{Akhoundi et al.\ 2025} --- fabricated 1024-channel 40-nm event-driven detection + SOM clustering at 74 nW/ch; very-high-channel custom-silicon competitor, no wireless.
\item \textbf{Chen et al.\ 2026} --- 128-channel fabricated on-chip sorter; state-of-the-art custom-silicon competitor, accuracy/power benchmark, no wireless.
\item \textbf{Jiang et al.\ 2026} --- 1024-channel fabricated on-chip sorter; custom-silicon competitor at extreme channel count, no AFE, no wireless.
\item \textbf{Lukito et al.\ 2026} --- 65-nm analog-CIM fabricated on-chip sorter that integrates the AFE; custom-silicon competitor, highest reported on-chip accuracy, no wireless.
\item \textbf{Shaeri \& Sodagar 2020} --- on-implant classification framework (feature selection + window discrimination); layout/synthesis only, no fabricated chip, no wireless.
\item \textbf{Valencia \& Alimohammad 2019} --- real-time parallel-OSort sorter, FPGA prototype + estimated 32-nm ASIC; on-device unsupervised-sorting competitor, no commodity hardware, no wireless.
\item \textbf{Han et al.\ 2025 (L-Sort)} --- localization-based on-device sorter exploiting multi-channel spatial features, FPGA + simulated 180-nm ASIC; relevant to the spatiotemporal feature-space axis, no wireless.
\item \textbf{Melo-Thomas et al.\ 2017 (TWS)} --- 4-channel wireless behavioral-recording headstage in the exact application domain, but raw streaming with no on-device sorting; illustrates the low-channel-count ceiling.
\end{enumerate}

\end{document}
